% RESEARCH_DRAFT_SHA256: 48058d54cb13d702ea2d2c641bda48fc2acd80b0249991713befaa2d96ecef73
% Layout and source-reference wrappers expanded for the web renderer.
\documentclass[11pt]{article}
\usepackage{amsmath,amssymb,mathtools}
\usepackage{fontspec,unicode-math}
\setmainfont{Latin Modern Roman}
\setsansfont{Latin Modern Sans}
\setmathfont{Latin Modern Math}
\usepackage{xurl,hyperref}
\newcommand{\sref}[2]{\hyperlink{source:#1:#2}{[S#2]}}
\newcommand{\cataloguescope}[1]{\paragraph{Mathematical question.}#1}
\begin{document}
\section*{Research-attempt provenance}
This standalone research attempt retains the complete original definition body below,
including its historical source and preparation notes. Those notes describe the
earlier catalogue preparation; the new attempt and its verification are identified
separately after that body. No original ID, catalogue status, or source statement is changed.

% UnsolvedMath ID 5500010; source catalogue code AMR-054-0010
\section{Simple Linear-Time Polygon Triangulation}\label{5500010}
{\small UnsolvedMath ID 5500010 · Computer Science}
\subsection{Definitions and mathematical statement}

\paragraph{Shared notation.}$\mathbb N=\{1,2,\ldots\}$, and $\mathbb Z,\mathbb Q,\mathbb R,\mathbb C$ denote the integers, rationals, reals and complex numbers. Any other conventions are specified locally.


A simple polygon has a non-self-intersecting polygonal boundary with $n$ vertices. A triangulation partitions its interior into triangles by noncrossing diagonals, using its vertices. Linear time means $O(n)$ running time. Simplicity of the algorithm is the qualitative criterion used in the source.
\paragraph{Statement recorded in the supplied source.}
Find a deterministic linear-time triangulation algorithm for arbitrary simple polygons that is significantly simpler than Chazelle\textquotesingle s. The related question recorded with it asks for a simple linear-time algorithm computing a shortest path within a simple polygon without first using a more complicated triangulation algorithm. \href{https://github.com/edemaine/topp/blob/main/Problems/P.000010}{[S2]}
\subsection{Short English statement}
Find a deterministic linear-time triangulation algorithm for arbitrary simple polygons that is significantly simpler than Chazelle\textquotesingle s. The related question recorded with it asks for a simple linear-time algorithm computing a shortest path within a simple polygon without first using a more complicated triangulation algorithm.
\subsection{Sources}
\begin{description}
\item[S1] ULAM.AI and the UnsolvedMath Contributors, UnsolvedMath: A Curated Collection of Open Mathematics Problems, record 5500010 (AMR-054-0010). CC BY 4.0. \url{https://huggingface.co/datasets/ulamai/UnsolvedMath}.
\item[S2] Erik D. Demaine, Joseph S. B. Mitchell and Joseph O\textquotesingle Rourke (editors), \emph{The Open Problems Project}, maintained original source, Problem 10, Statement and complete Partial/Related Results. \url{https://github.com/edemaine/topp/blob/main/Problems/P.000010}.
\end{description}
\label{end:5500010}

\section{Literature and context check}
The maintained TOPP Problem 10 page was consulted. It retains both the
simpler deterministic triangulation question and the related shortest-path
question. Its discussion distinguishes randomized constructions from the requested
deterministic guarantee. ``Significantly simpler'' has no formal measure in the
source. No recent-resolution claim is made from that historical page alone.
\section{Research attempt: Auditing fan triangulation and ear-scan costs}
The first attempted linear algorithm was to choose one boundary vertex and
draw a fan. For a convex polygon $v_0,\ldots,v_{n-1}$ this works: each diagonal
$v_0v_i$ lies in the polygon by convexity, and the triangles
$(v_0,v_i,v_{i+1})$, $1\le i\le n-2$, have disjoint interiors and exhaust it.
One proof is to intersect every ray from $v_0$ with the convex polygon: its
intersection is an interval, whose outer endpoint lies on exactly one opposite
edge except at finitely many directions. Thus construction and explicit output
cost $O(n)$ arithmetic operations. For two points of a convex polygon, their
straight segment is also a shortest internal path: every rectifiable competitor
has length at least the Euclidean distance, by applying the triangle inequality
to its polygonal approximations.

The fan inference fails on this explicit simple polygon, in boundary order:
\[
 (0,0),(4,0),(4,4),(3,4),(3,1),(1,1),(1,4),(0,4).
\]
Its interior is the outer rectangle with the upper open rectangular notch
$1<x<3$, $1<y<4$ removed. The proposed diagonal from $(0,0)$ to $(4,4)$
passes through $(2,2)$, outside the polygon. Its boundary edges have disjoint
relative interiors, so this is a valid simple-polygon counterexample to the fan
algorithm, rather than a self-intersecting input. A shortest-path shortcut through
the notch likewise fails the domain constraint.

The next attempt was ear removal with an explicit emptiness scan: when $m$
vertices remain, select a convex candidate, test all other $m-3$ vertices for
membership in its triangle, and remove it if empty. On a strictly convex polygon
every such candidate succeeds, but the specified verifier still makes exactly
\[
 \sum_{m=4}^n(m-3)=\frac{(n-3)(n-2)}2
\]
point-in-triangle tests. The input family can be the vertices of any strictly
convex $n$-gon. Thus constant amortized time per removed ear has not been proved:
even a sequence of successful ears can require quadratic verification under the
chosen implementation. This is a running-time counterexample to that concrete
implementation, not a lower bound for ear methods with more careful maintenance.

The two failures locate the remaining operation: cheaply certifying that a
candidate diagonal is internal in arbitrary nonconvex geometry. Neither the
convex fan proof nor the successful-ear count pays for that operation. A
replacement visibility structure would require its own preprocessing and update
analysis before an $O(n)$ total bound could be asserted.
\section{Verification and adversarial audit}
The notch example was checked by its coordinate intervals and by the point
$(2,2)$ on the proposed diagonal. The exact quadratic sum was evaluated
symbolically. Triangles ($n=3$) give zero ear-removal tests, consistently. Collinear
boundary vertices and shortest paths in nonconvex polygons are not claimed
handled by the convex special case.
\paragraph{Audit corrections.} The quadratic count is only for the specified full-scan verifier, not a general lower bound; convexity was not silently imposed on the original target.
\paragraph{External computation.} None; symbolic arguments and boundary checks only.
\section{Scope and remaining gaps}
No new algorithm for arbitrary simple polygons and no general internal
shortest-path algorithm was obtained. The first unresolved step is a simple
deterministic linear-total-cost mechanism for internal-diagonal certification;
both bundled questions and their qualitative simplicity requirement remain.
\paragraph{Integration review.} Codex (GPT-6; exact serving revision not exposed) reviewed the written arguments and expanded layout and source-reference wrappers. Repository metadata and the production web renderer were checked using Python 3.13 and JavaScript with MathJax 3.2.2. This model-assisted review is not independent human proof verification; no PDF compilation or formal proof checking is asserted.

\section{Completion Estimate}
\noindent\textbf{COMPLETION: 0\%}
The calculations below the preserved source statement investigate a concrete route,
but establish no substantiated new progress toward the entire original target.
Known special cases and failed sufficient conditions are not credited as a solution.
\section{References}
The original source references and catalogue attribution are preserved in the
Sources section above. The following URLs were actually consulted for this attempt;
retention of an additional original reference does not claim fresh consultation of it.
\begin{itemize}
\item \url{https://topp.openproblem.net/p10} (consulted 2026-10-11). 
\end{itemize}
\section{Model and contributor disclosure}
OpenAI gpt-6-astra, requested reasoning configuration ultra, generated this research
attempt on 2026-10-11. The exact serving revision is not exposed. Codex GPT-6
performed orchestration and repository-format review; its exact serving revision
is also not exposed. The mathematical route, boundary checks and final adversarial
self-audit followed contribution prompts A then B in
\texttt{docs/contribution\_prompts.md}, after reading \texttt{CONTRIBUTING.md}.
Python 3.14.2 assembled the standalone documents; any mathematical computation is
specified in Verification. No theorem prover was used. No human mathematical
authorship, substantial human proof correction, or independent human proof checking
is asserted. Model-assisted review, source retrieval, and structural checks do not
establish correctness or novelty. Full private reasoning transcripts are unavailable
as reproducibility artifacts; the complete arguments, inputs and outcomes actually
used are stated here or in the companion batch audit.
\section{Final status}
Unresolved. This is a research attempt with an explicit failed route or restricted
calculation, not a full solution, not a current-status certification, and not a claim
that the elementary material is new. The catalogue's original source-status label
is preserved independently of this attempt status.
\end{document}
